\bookchapter{第四十五章：质量保障——让AI写的代码可靠}

\chapterintro{AI 写的代码有一个特点：它通常“看起来很对”——变量命名规范、注释齐全、逻辑清晰。但“看起来对”不等于“真的对”。AI 代码中隐藏的问题往往比手写代码更难发现，因为它的表面质量太高了，容易让人放松警惕。本章教你如何为 AI 代码建立有效的质量保障体系——用测试、审查、版本控制三道防线，把“看起来对”变成“真的对”。}

\booksection{一、测试驱动：让 AI 写测试}

测试是 AI 编程质量保障的第一道防线。好消息是：\textbf{让 AI 写测试比让它写功能代码更容易、更可靠}——因为测试的预期行为是明确的，不太容易产生“看起来对但实际不对”的代码。

测试驱动的 AI 编程流程是：先让 AI 写功能代码，验证基本功能通过后，立即让它写测试用例。测试用例应覆盖正常路径、边界情况和异常情况。然后再让 AI 根据测试反馈修正功能代码。这个流程的好处是：测试一旦写好，就是持续的验证工具——后续任何修改都可以通过运行测试来确认没有引入回归问题。

\begin{lstlisting}
请为 calculate_discount 函数编写单元测试。
测试用例应覆盖：
1. 正常路径：
   - 100元商品打8折 -> 80
   - 200元商品打5折 -> 100
2. 边界情况：
   - 0元商品打8折 -> 0
   - 100元商品打0折 -> 100（不打折）
   - 100元商品打10折 -> 100（全价）
3. 异常情况：
   - 负数价格 -> 应抛出错误
   - 折扣率 > 1 -> 应抛出错误
   - 非数字输入 -> 应抛出错误
使用Vitest框架，给出完整测试文件。
\end{lstlisting}

注意这个指令的特点：它\textbf{明确列出了三类测试用例}（正常、边界、异常），不给 AI 自由发挥的空间去“选择性测试”。AI 有时候会只写正常路径的测试来“快速交差”，明确列出所有类别可以避免这个问题。

\booksection{二、代码审查技巧}

即使有测试覆盖，你仍然需要对 AI 产出的代码进行人工审查。代码审查不是逐行阅读——那太耗时也太依赖你的编程水平。有效的 AI 代码审查应该关注以下几个关键维度：

\begin{center}
\begin{tabularx}{\linewidth}{|X|X|X|}
\hline
\textbf{审查维度} & \textbf{关注什么} & \textbf{常见问题} \tabularnewline
\hline
安全性 & 输入验证、密钥处理、SQL注入防护 & AI经常忘记验证输入，密钥硬编码在代码中 \tabularnewline
错误处理 & 异常是否被捕获，错误信息是否合理 & AI喜欢用空catch块或只console.log \tabularnewline
性能 & 是否有明显的性能问题，如N+1查询、无限循环 & AI有时会写出O(n²)的算法 \tabularnewline
可维护性 & 代码结构是否清晰，命名是否有意义 & AI偶尔会生成过度抽象的代码 \tabularnewline
依赖管理 & 引入的依赖是否必要，版本是否明确 & AI倾向于引入不必要的库 \tabularnewline
\hline
\end{tabularx}
\end{center}

\begin{center}\small\color{gray}AI 代码审查的关键维度\end{center}

审查时不必自己逐一检查——可以让 AI 帮你做第一轮审查。指令如：“请审查这段代码的安全性、错误处理和性能，指出潜在问题。”AI 的自我审查能力在不断提升，虽然不能完全替代人工审查，但能帮你过滤掉大部分明显问题。

\booksection{三、版本控制：Git 在 AI 编程中的角色}

Git 在 AI 编程中的角色比在传统编程中更加重要。因为 AI 的每次修改都可能引入不可预知的变化，你需要一个可靠的“后悔药”——\textbf{版本控制就是那个后悔药}。

AI 编程中 Git 的最佳实践是“频繁提交、小步快进”。每完成一个可验证的功能点，就提交一次。这样当 AI 的某次修改出了问题时，你可以精准回退到上一个可用状态，而不需要推翻所有工作。

另一个实用技巧是“分支隔离”。对于不确定结果的 AI 修改（比如大步重构），先创建一个新分支来做。如果结果满意，合并回主分支；如果不满意，直接删除分支，不影响主分支。这比在主分支上反复修改安全得多。

\begin{enumerate}
  \item 每次 AI 修改前，确保工作区是干净的（没有未提交的更改）。这样如果 AI 的修改不满意，可以用\texttt{git checkout .}快速回退。
  \item 每次 AI 修改后、验证通过后，立即提交。提交信息描述这次修改做了什么，比如"feat: 添加用户注册功能"。
  \item 对于大步重构，先创建分支：\texttt{git checkout -b refactor/auth-system}。在分支上完成修改，测试通过后再合并。
  \item 养成查看 diff 的习惯。每次 AI 修改后，用\texttt{git diff}查看具体改了什么。这不仅能帮你发现问题，还能帮你学习代码。
\end{enumerate}

\booksection{四、常见陷阱与避坑指南}

AI 编程中有几个反复出现的陷阱，几乎所有初学者都会踩。提前知道它们，能帮你少走很多弯路。

\booksubsection{陷阱一：信任陷阱}

“AI 写的代码看起来很专业，一定没问题吧？”——这是最危险的陷阱。AI 生成的代码在表面上非常整洁——命名规范、注释齐全、结构清晰——这种“表面专业感”会让人放松警惕。但\textbf{表面整洁不代表逻辑正确}。\textbf{解药}：永远通过测试验证，而非通过“看起来对不对”判断。

\booksubsection{陷阱二：膨胀陷阱}

“既然 AI 能帮我写代码，那不如把功能做得全面一点？”——这会导致代码膨胀。AI 很擅长“多给”：你让它做一个登录页面，它会给你一套完整的用户管理系统。多余的功能不仅增加复杂度，还增加了出 bug 的概率。\textbf{解药}：严格遵循 Simplicity First 原则，明确告诉 AI“只做这个功能，不要额外添加”。

\booksubsection{陷阱三：理解陷阱}

“AI 应该理解我的意思吧？”——不一定。AI 的理解能力虽然很强，但对模糊的需求仍然会“猜”。当它猜错了，产出的代码就是错的。更麻烦的是，它猜错的代码通常“看起来很合理”，让你误以为它理解了。\textbf{解药}：在指令中提供足够的具体细节，减少 AI 需要“猜”的空间。

\booksubsection{陷阱四：上下文陷阱}

“我在之前的对话里说过了，AI 应该还记得吧？”——很可能不记得。随着对话变长或开新会话，AI 会丢失早期上下文。如果某个关键约束没有被写入项目记忆文件，它很可能在后续修改中被忽略。\textbf{解药}：关键约束写入项目记忆文件，不要依赖对话记忆。

\begin{tcbwarn}{四大陷阱速记}
信任陷阱——别信“看起来对”；膨胀陷阱——AI 爱“多给”；理解陷阱——模糊需求它会“猜”；上下文陷阱——对话记忆会“忘”。对应的解药分别是：测试验证、砍掉多余、给足细节、写入记忆文件。

\end{tcbwarn}

\begin{tcbpractice}{来动手：给你的项目建三道防线}
如果你已经开始做项目一（工具站）：(1) 让 AI 给核心函数写 3 类测试用例；(2) 按五个维度审查一遍代码（安全性/错误处理/性能/可维护性/依赖）；(3) 用 git 初始化项目，做第一次提交。三道防线齐了，你的项目才真正属于你。

\end{tcbpractice}

\begin{tcbsummary}{本章小结}
\begin{itemize}
  \item AI 代码"看起来对"$\neq$"真的对"——表面质量太高反而容易放松警惕
  \item 让 AI 写测试：明确列出正常/边界/异常三类用例，防止"选择性测试"
  \item 审查五维度：安全性、错误处理、性能、可维护性、依赖管理
  \item Git 是"后悔药"：频繁提交、小步快进、分支隔离、常看 diff
  \item 四大陷阱：信任（测试验证）、膨胀（砍掉多余）、理解（给足细节）、上下文（写入记忆文件）
\end{itemize}

\end{tcbsummary}
